\section{Conclusion}

We presented a systematic comparison of Rule-Based SAST tools, general-purpose
LLM scanners, and Security-Specialized systems, evaluated on
\textsc{RealVuln}, the first fully open-source benchmark built from real-world
vulnerable code.  Version~1 spans 26 repositories, 796 manually labeled
findings, and 120 false-positive traps across 15 scanners in three categories.
Every artifact (ground truth, raw scanner output, scoring pipeline, dashboard
generator, and evaluation prompts) is released under a permissive license
with a reproducible manifest that version-locks results to exact benchmark
snapshots.

Three tiers emerge clearly. The Security-Specialized scanner Kolega.Dev (F3\,=\,73.0 strict) leads, showing that purpose-built security architecture delivers higher recall than general-purpose methods. General-purpose LLM scanners form a competitive middle tier. The best, Claude Sonnet~4.6 (F3\,=\,51.7), scores nearly 3$\times$ higher than the best Rule-Based SAST tool, Semgrep (F3\,=\,17.7). No scanner dominates across all CWE families. Rule-Based tools retain a precision advantage on well-specified injection patterns. LLM-based systems do better on logic-level and authentication flaws that need broader code context. The contrast between the two Security-Specialized systems is instructive. The SecLab Taskflow Agent uses sophisticated multi-stage threat modeling with GPT-5.4 but scores F3\,=\,8.4, below every Rule-Based SAST tool. Its depth-first design audits only 3--5 threat categories per repository. Breadth of coverage, not just Security-Specialized architecture, separates the top tier from the rest. The future of vulnerability detection is not in applying general-purpose AI to security prompts. It is in building Security-Specialized systems that combine LLM reasoning with domain-specific architecture \emph{and} broad detection coverage.

The benchmark is designed to grow. New repositories, scanners, and ground-truth corrections can be added without invalidating prior results (\S\ref{sec:living-benchmark}).  Version 2 will expand coverage to JavaScript/TypeScript, Go, and
Java, and will introduce Type 2 targets drawn from real production CVEs,
setting a higher bar for precision in realistic, low-vulnerability-density
codebases.  We invite the security research community to submit scanner results,
propose new repositories, and challenge existing labels.

All code, data, results, and tooling are available at \url{https://github.com/kolega-ai/RealVulnBenchmark}. A live interactive dashboard is hosted at \url{https://realvuln.kolega.dev}. We encourage the security community to verify our results, challenge our ground truth, and submit their own scanner evaluations.
